\section{Kodowanie, diagonalizacja i prawda: sens twierdzenia Tarskiego}\label{sec:sens-tarski}
\subsection*{Po co zamieniać napisy i obliczenia na liczby}\label{sec:18:1}
Numeracja Gödla nie nadaje liczbom tajemniczych znaczeń. Jest sposobem zapisu danych, podobnie jak kodowanie tekstu w komputerze: skończony napis lub skończone obliczenie otrzymuje jednoznaczny numer. Reprezentowalność mówi następnie, że poprawność konkretnych operacji na takich danych można wyrazić i sprawdzić w arytmetyce.

Dzięki temu arytmetyka potrafi mówić nie tylko o dodawaniu liczb, ale także o tym, że dana liczba koduje formułę, podstawienie albo dowód. To jest techniczny powód, dla którego teoria liczb może opisywać własną działalność dowodową. Lemat diagonalny pozwala z takiego opisu zbudować zdanie dowodliwie równoważne stwierdzeniu o jego własnym kodzie. Nie trzeba dodawać do języka magicznego wyrażenia „to właśnie zdanie”.

Ostrożność nadal jest potrzebna: reprezentowalność obliczeń na każdym konkretnym wejściu nie jest automatycznie dowodem, że program kończy się na wszystkich wejściach. Dla funkcji p.r. uzasadniliśmy odpowiednią totalność w PA w Lemacie~\ref{res:11:3}; nie zrobiliśmy tego dla każdej totalnej funkcji obliczalnej. Arytmetyzacja nie zmienia nieograniczonego szukania w ograniczone sprawdzanie gotowego świadectwa.

\subsection*{Dlaczego zdanie Gödla nie jest zwykłym paradoksem kłamcy}\label{sec:18:2}
Zdanie „to zdanie jest fałszywe” próbuje użyć własnego predykatu prawdy bez ograniczenia jego zakresu. Zdanie Gödla mówi zamiast tego o niedowodliwości w konkretnej teorii. Te pojęcia nie są zamienne: prawdziwe zdanie może nie mieć dowodu z wybranego zbioru aksjomatów. Właśnie ta różnica pozwala uzyskać niezupełność, a nie sprzeczność.

Twierdzenie Tarskiego pokazuje z kolei, co poszłoby źle, gdyby arytmetyka miała jedną formułę $\operatorname{Tr}(x)$ poprawnie rozpoznającą prawdę wszystkich własnych zdań. Diagonalizacja zastosowana do $\neg\operatorname{Tr}(x)$ dawałaby już prawdziwy konflikt typu kłamcy. Wniosek brzmi: taka formuła nie istnieje, a nie: arytmetyka zawiera sprzeczność.

\subsection*{Prawda istniejąca z zewnątrz a prawda zdefiniowana wewnątrz}\label{sec:18:3}
Twierdzenie Tarskiego nie zabrania mówić o prawdzie w $\mathbb N$ w bogatszym języku metateorii, na przykład teorii mnogości. Tam możemy opisać strukturę $\mathbb N$, wartościowania i rekurencyjną definicję spełniania. Nie zabrania też rozpatrywania samego zbioru $\operatorname{Th}(\mathbb N)$ z zewnątrz. Zabrania zdefiniowania pełnego zbioru prawd arytmetycznych za pomocą jednej formuły tego samego języka arytmetyki.

Można również definiować prawdę dla odpowiednio ograniczonych klas formuł. Zakaz dotyczy jednego predykatu obejmującego wszystkie zdania własnego języka, a nie każdego lokalnego testu prawdziwości. W wersji syntaktycznej żadna niesprzeczna $T\supseteq Q$ nie może udowodnić wszystkich równoważności $\varphi\leftrightarrow\operatorname{Tr}(\ulcorner\varphi\urcorner)$ przy jednym arytmetycznym $\operatorname{Tr}$.

\subsection*{Trzy różne bariery: definiowanie, wyliczanie, rozstrzyganie}\label{sec:18:4}
Warto oddzielić wyniki, które na pierwszy rzut oka brzmią podobnie:
\begin{itemize}
\item \emph{Definiowalność:} czy jedna formuła opisuje dany zbiór? Dowodliwość $\operatorname{Bew}_T$ jest definiowalna w arytmetyce; pełna prawda arytmetyczna nie jest, na mocy twierdzenia Tarskiego.
\item \emph{Efektywne wyliczanie:} czy program może wypisać wszystkie elementy zbioru, nie wypisując innych? Twierdzenia efektywnej teorii można wyliczać, ale wszystkich prawd $\operatorname{Th}(\mathbb N)$ nie można, na mocy Wniosku~\ref{res:9:5}.
\item \emph{Rozstrzyganie:} czy program zawsze kończy pracę z poprawną odpowiedzią „należy” lub „nie należy”? Dla zbioru twierdzeń rozpatrywanych niesprzecznych rozszerzeń $Q$ jest to niemożliwe, mimo możliwości ich wyliczania.
\end{itemize}
Brak efektywnego wyliczania nie oznacza sam przez się braku definicji: formuły mogą opisywać zbiory niebędące rekurencyjnie przeliczalnymi. Z kolei przeliczalność nie oznacza efektywnej wyliczalności. $\operatorname{Th}(\mathbb N)$ jest podzbiorem przeliczalnego zbioru zdań, więc jest co najwyżej przeliczalny, ale żadna maszyna nie wypisze dokładnie wszystkich jego elementów. To kolejny powód, dla którego teoriomnogościowego wyboru z dowodu Löwenheima--Skolema nie należy mylić z algorytmem.

